%% ===========================================================================
\section{Coefficients of the rectangular, elliptical and ellipsoidal averaging
windows}
\label{app:Mcoeff}
%% ===========================================================================

For a uniform kernel over a window $\mathcal{W}$ centred at the origin, the
normalised second moment is
\begin{equation}
  M_{mn} = \frac{1}{|\mathcal{W}|}\int_{\mathcal W} s_m s_n \,\mathrm{d}\bm{s},
  \label{eq:appM}
\end{equation}
which is diagonal for a window whose axes are aligned with the coordinate axes
and which transforms as $M\mapsto RMR^{\mathsf T}$ under a rotation $R$ of the
window. The three closed forms used in the paper follow from
\eqref{eq:appM} by elementary integration.

\paragraph{Cuboid.} For $\mathcal W=[-a_1,a_1]\times[-a_2,a_2]$ the area is
$4a_1a_2$ and
\begin{equation}
  M_{11} = \frac{1}{4a_1a_2}\int_{-a_1}^{a_1}\!\!\int_{-a_2}^{a_2}
           s_1^{2}\,\mathrm{d}s_1\,\mathrm{d}s_2
         = \frac{1}{4a_1a_2}\cdot 2a_2\cdot\frac{2a_1^{3}}{3}
         = \frac{a_1^{2}}{3},
  \label{eq:Mcub}
\end{equation}
and identically $M_{22}=a_2^{2}/3$; the off-diagonal integral vanishes by the
oddness of $s_1s_2$. A three-dimensional cuboid of half-widths $a_i$ gives the
same coefficient along each axis.

\paragraph{Ellipse.} For $\mathcal W=\{s: (s_1/a_1)^{2}+(s_2/a_2)^{2}\le1\}$ the
area is $\pi a_1a_2$ and the polar second moment of the semi-axis is
$\int_{\mathcal W}s_1^{2}\,\mathrm{d}A = \pi a_1^{3}a_2/4$, so that
\begin{equation}
  M_{11} = \frac{\pi a_1^{3}a_2/4}{\pi a_1a_2} = \frac{a_1^{2}}{4},
  \qquad M_{22}=\frac{a_2^{2}}{4}.
  \label{eq:Mell}
\end{equation}

\paragraph{Ellipsoid.} For the three-dimensional ellipsoid of semi-axes $a_i$
the volume is $4\pi a_1a_2a_3/3$ and
$\int_{\mathcal W}s_1^{2}\,\mathrm{d}V = 4\pi a_1^{3}a_2a_3/15$, hence
\begin{equation}
  M_{11} = \frac{a_1^{2}}{5},\qquad
  M_{22}=\frac{a_2^{2}}{5},\qquad
  M_{33}=\frac{a_3^{2}}{5}.
  \label{eq:Mellip}
\end{equation}

\paragraph{Equal-second-moment windows.} Equations~\eqref{eq:Mcub}--\eqref{eq:Mellip}
are the reason the shape of the window drops out of the model. A cuboid of
half-width $a$ and an ellipsoid of semi-axis $l$ carry the same second moment
when $a^{2}/3=l^{2}/5$, that is when $a=l\sqrt{3/5}$; for an ellipse the same
condition reads $a=l\sqrt{3/4}$. Two windows related in this way produce the
same gradient stiffness $K_g$ and therefore the same solution, which is what
Fig.~\ref{fig:domain} and the polar panel of Fig.~\ref{fig:petals} show
numerically, and what the ratio $5/3$ between the coefficients of an equal
half-width pair expresses analytically.

\paragraph{Rotated window.} A rectangular window whose long axis makes an angle
$\theta$ with $x$ has
\begin{equation}
  M_{11}(\theta) = M_{11}\cos^{2}\theta + M_{22}\sin^{2}\theta,\qquad
  M_{22}(\theta) = M_{11}\sin^{2}\theta + M_{22}\cos^{2}\theta,\qquad
  M_{12}(\theta) = \tfrac12\left(M_{22}-M_{11}\right)\sin2\theta,
  \label{eq:Mrot}
\end{equation}
so that the off-diagonal coupling is largest at $\theta=45^\circ$ and vanishes
when the window is aligned with the axes. This is the tensor that the
orientation sweep of Sec.~\ref{sec:orientation} rotates.

\paragraph{Reduction to a scalar length.} The ellipsoid gives the coefficient
that the one-scale theory of the authors' earlier work carries:
$M_{mn}=L_{mn}/5$, hence $\tfrac14M_{mn}=\tfrac1{20}L_{mn}$ and the gradient
energy of~\cite{fem1} is recovered exactly, and setting
$M_{mn}=\ell^{2}\delta_{mn}$ reduces the same expression to the isotropic
Aifantis form, for which the boundary-layer length of a clamped bar is
$\sqrt{\ell^{2}/2}$, the value that follows from putting $L_{mn}=5\ell^{2}
\delta_{mn}$ in the $\sqrt{L/10}$ of the one-scale theory.
